求 \(f(x)=0\) 时，我们只能查询函数值，有时还能得到导数。取 \(f(x)=x^3-2x-5\)，它的唯一实根在 2 和 3 之间。二分法每次折半区间，约 10 步多 3 位十进制精度，不问 \(f\) 长什么样。Newton 法从 \(x_0=2\) 出发，第一步就到 \(x_1\approx2.1\)，第二步 \(x_2\approx2.0946\)，两步就拿到二分法需要 15 步的精度——前提是初值足够近。

二分法用一个含根区间担保收敛；Newton 法用局部导数换取二次速度；割线法则用历史函数值近似导数。每种方法维护不同的证据，进入快收敛区的条件不同，停止条件能证明的事也不同。

\begin{enumerate}
\tightlist
\item
  括区与二分法
\item
  不动点与 Newton 法
\item
  割线法、收敛阶与混合策略
\end{enumerate}

三种方法可以看成三种姿态。二分法始终保存"根还在这个区间里"这一事实；Newton 法相信当前切线在附近足够像原函数；割线法用两个旧函数值猜斜率。方法快不快，要连它放弃了什么安全网一起看。

前一单元：\hyperref[ch:061]{``浮点数与误差：计算机为什么不等于实数系统''}。

\subsection{本章篇目}

以下篇目按概念依赖排列；若从具体问题进入，也可以先打开最接近当前困惑的一篇，再沿篇内链接回补。

\begin{itemize}
\tightlist
\item \hyperref[sec:08-Numerical-Analysis/02-Root-Finding/01]{``括区与二分法''}；
\item \hyperref[sec:08-Numerical-Analysis/02-Root-Finding/02]{``不动点与 Newton 法''}；
\item \hyperref[sec:08-Numerical-Analysis/02-Root-Finding/03]{``割线法、收敛阶与混合策略''}。
\end{itemize}
